% !TEX program = xelatex
%
% Title page — submitted as a separate item, not part of the anonymised
% manuscript. This layout is correct whether JSS reviews single-anonymised
% (author names visible to reviewers) or double-anonymised (they are not):
% Elsevier's submission system has a separate "Title Page" item either way.
%
% FILL IN: the three fields marked <<<  >>> below, then delete this comment
% block. Everything else is already worded for submission.

\documentclass[preprint,12pt]{elsarticle}

\usepackage{graphicx}
\usepackage{booktabs}
\usepackage{amsmath,amssymb}
\usepackage{siunitx}
\usepackage[hidelinks]{hyperref}
\usepackage{xcolor}

\journal{Journal of Systems and Software}

\begin{document}

\begin{frontmatter}

\title{Identifiability and Confounding in Test-Adequacy Metric Validity
Studies: A Cross-Language Empirical Analysis}

\author[a]{<<< AUTHOR NAME >>>}
\affiliation[a]{<<< AFFILIATION >>>}

\ead{<<< CORRESPONDING E-MAIL >>>}

\end{frontmatter}

\section*{CRediT authorship contribution statement}

\textbf{<<< AUTHOR NAME >>>:} Conceptualization, Methodology, Software,
Validation, Formal analysis, Investigation, Data curation, Writing -- original
draft, Writing -- review \& editing, Visualization.

\section*{Declaration of competing interest}

The author declares no known competing financial interests or personal
relationships that could have appeared to influence the work reported in this
paper.

\section*{Funding}

This research did not receive any specific grant from funding agencies in the
public, commercial, or not-for-profit sectors.

\section*{Acknowledgements}

The author thanks the maintainers of Defects4J and BugsInPy for making the
benchmarks publicly available, and the ClawsGO Science Harness for the
computing environment used for the measurement runs.

\section*{Data availability}

All scripts, collected measurements, and analysis code are available at
\url{https://github.com/JunjieChen0/test-adequacy-resolution}. Every number
reported in the paper is produced by \texttt{scripts/analyze\_all.py} from the
shipped measurements.

\end{document}
